\documentclass[11pt]{article}

% ===================== PACKAGES =====================
\usepackage[a4paper,margin=1in]{geometry}
\usepackage{amsmath,amssymb}
\usepackage{enumitem}
\usepackage{fancyhdr}
\usepackage{xcolor}

% ===================== HEADER & FOOTER =====================
\pagestyle{fancy}
\fancyhf{}
\lhead{CSE 400: Fundamentals of Probability in Computing}
\rhead{Lecture 21}
\cfoot{\thepage}

% ===================== TITLE =====================
\title{
    \normalsize School of Engineering and Applied Science (SEAS), Ahmedabad University \\
    \vspace{0.2cm}
    \textbf{CSE 400: Fundamentals of Probability in Computing}\\
    \Large Lecture Scribe Submission
}
\author{}
\date{}

\begin{document}
\maketitle

\vspace{-2cm}
\begin{center}
    \begin{tabular}{ll}
        \textbf{Name:} & Ishika Patel \\[1.5ex]
        \textbf{Enrollment No:} & AU2440088 \\[1.5ex]
        \textbf{Group:} & s1\_g1\_fin \\[1.5ex]
        \textbf{Lecture:} & 21 \\[1.5ex]
        \textbf{Instructor:} & Dhaval Patel, PhD \\[1.5ex]
        \textbf{Date:} & March 30, 2026
    \end{tabular}
\end{center}

\hrule
\vspace{0.5cm}


\section*{Lecture 21: Tail Probabilities and Bounds}

\section*{1. Tails}

\textbf{Definition}

The tail of a random variable $X$ is:
\[
P(X > x)
\]

\textbf{Why do we care about tails?}

\begin{itemize}
\item Jobs delayed more than 24 hours
\item Packets dropped due to full buffer
\end{itemize}

\textbf{Key idea}

\begin{itemize}
\item Mean $\rightarrow$ average behavior (easy)
\item Tail $\rightarrow$ rare extreme events (hard)
\end{itemize}

\textbf{Why are tails hard to compute?}

Requires summing many small probabilities

\section*{2. Tails Example 1/2}

\textbf{Question}

Suppose we’re distributing $n$ jobs among $n$ servers at random. What’s the probability that a particular server gets $\geq k$ jobs?

\textbf{Intuition}

\begin{itemize}
\item Jobs are assigned randomly
\item Most servers get few jobs
\item Occasionally, one server gets many jobs
\item This is a tail event (rare overload)
\end{itemize}

\textbf{Model}

\[
X \sim \text{Binomial}(n, p)
\]

\textbf{Probability}

\[
P(X \geq k) = \sum_{i=k}^{n} \binom{n}{i} p^i (1-p)^{n-i}
\]

\textbf{Note}

Highlighted server has $\geq k$ jobs

\section*{3. Tails Example 2/2}

\textbf{Question}

Jobs arrive to a datacenter according to a Poisson process with rate $\lambda$ jobs/hour. What’s the probability that $\geq k$ jobs arrive during the first hour?

\textbf{Intuition}

\begin{itemize}
\item Jobs arrive randomly over time
\item Most hours have moderate arrivals
\item Sometimes, we see a burst of arrivals
\item This is a tail event (sudden spike in load)
\end{itemize}

\textbf{Model}

\[
X \sim \text{Poisson}(\lambda)
\]

\textbf{Probability}

\[
P(X \geq k) = \sum_{i=k}^{\infty} e^{-\lambda} \frac{\lambda^i}{i!}
\]

\textbf{Note}

Many arrivals within one hour ($\geq k$ jobs)

\section*{4. Why Tail Bounds?}

\textbf{Observation}

Exact tail probabilities are often hard to compute

\textbf{Exact computation}

Complicated, long sums

\textbf{Idea: Instead of exact value}

\begin{itemize}
\item We find an upper bound
\item Easier to compute
\item Still gives a guarantee
\end{itemize}

\textbf{Form}

\begin{itemize}
\item Exact:
\[
P(X \geq k)
\]
\item Bound:
\[
P(X \geq k) \leq \text{something simple}
\]
\end{itemize}

\textbf{Key idea}

Approximate safely instead of computing exactly

\section*{5. Running Example (Key Idea)}

\textbf{Goal}

Develop progressively better (tighter) tail bounds

Test all bounds on the same example

\textbf{Experiment}

Flip a fair coin $n$ times

\textbf{Distribution}

\[
X = \text{number of heads}
\]
\[
X \sim \text{Binomial}(n, \tfrac{1}{2})
\]

\textbf{Mean}

\[
E[X] = \frac{n}{2}
\]

\textbf{Question}

\[
P\left(X \geq \frac{3n}{4}\right)
\]

\textbf{Observation}

Tail region: getting unusually many heads

\section*{6. Markov’s Inequality}

\textbf{Key Idea}

Uses only the mean

Large values $\Rightarrow$ rare

\[
P(X \geq k) \leq \text{something simple}
\]

\textbf{Example}

Avg = 10 $\Rightarrow$ few values can be 100+

\textbf{Markov Inequality}

\[
P(X \geq a) \leq \frac{E[X]}{a}
\]

\textbf{Proof (Idea)}

\[
E[X] = \sum_{x=0}^{\infty} x \cdot p_X(x)
\]

\[
\geq \sum_{x=a}^{\infty} x \cdot p_X(x)
\]

\[
\geq \sum_{x=a}^{\infty} a \cdot p_X(x)
\]

\[
= a \cdot P(X \geq a)
\]

\[
\Rightarrow P(X \geq a) \leq \frac{E[X]}{a}
\]

\bigskip
\hrule
\vspace{0.2cm}
\begin{center}
    \small \textit{End of Lecture 21}
\end{center}

\end{document}